\documentclass[11pt]{amsart}
\usepackage{amssymb,amsmath}
\usepackage[hidelinks]{hyperref}
\newtheorem{theorem}{Theorem}
\newtheorem{lemma}[theorem]{Lemma}
\setlength{\emergencystretch}{2em}
\addtolength{\textheight}{3\baselineskip}
\hypersetup{pdftitle={Extreme zeros of ordered set-partition polynomials at a fixed block-size threshold},pdfauthor={Bowen Zuo}}
\title[Extreme zeros of ordered set partitions]{Extreme zeros of ordered set-partition polynomials at a fixed block-size threshold}
\author{Bowen Zuo}
\address{East China University of Science and Technology}
\date{November 27, 2026}
\begin{document}
\begin{abstract}
Give each block of size at least $j$ in a labelled ordered set partition
weight $t$. For every fixed $j\ge2$, any fixed number of zeros of
largest modulus are simple and negative in all sufficiently large rows.
These outer prefixes alternate pairwise throughout a window of $j+1$
consecutive rows, with orientation determined by the degrees. The
reciprocal polynomial has a Mittag--Leffler limit. Its first correction
orders corresponding roots one full threshold apart, where the leading
terms coincide.
\end{abstract}
\maketitle

\section{Statement}
Let $O_{n,j}(t)$ sum $t^{\#\{i:|B_i|\ge j\}}$ over lists
$(B_1,\ldots,B_\ell)$ of disjoint nonempty blocks with union $[n]$,
so each ordering of an unordered partition with $\ell$ blocks is counted.
The empty list has weight one. The degree is $\lfloor n/j\rfloor$,
attained by taking blocks of size $j$ and the remaining singletons.
The all-singleton lists give a positive constant term. Alexandersson
conjectures real-rootedness and consecutive interlacing for these
ordered rows~\cite{catalog}. We treat a fixed number of roots while
the degree increases. For the separate unordered model at threshold four,
Fang~\cite[Theorems 1 and 4]{fang} proves eventual interlacing of the
whole rows, using the same limiting functions up to constant factors
and an outer-root correction at distance four on the $d^5$ scale.
For ordered rows the scale is $d^j$, and the two-term expansion below
holds at every fixed threshold. Put
$b=1+1/(j+1)$, $G=j!b^j$ and
\[
 F_r(w)=\sum_{h\ge0}\frac{w^h}{(jh+r)!},\qquad 0\le r<j .
\]
Write $-\lambda_{r,k}$ for its zeros in increasing order of modulus,
starting at $k=1$.

\begin{theorem}\label{thm}
Fix integers $j\ge2$ and $K\ge1$. For all sufficiently large $n$,
the $K+1$ zeros of greatest modulus of $O_{n,j}$ are simple and negative.
Write them as $x_{n,1}<\cdots<x_{n,K+1}<0$. Every remaining
zero has modulus strictly less than $|x_{n,K+1}|$.
For $n<m\le n+j$ and $1\le k\le K$,
\[
\begin{array}{ll}
 x_{n,k}<x_{m,k}<x_{n,k+1},&\lfloor n/j\rfloor=\lfloor m/j\rfloor,\\
 x_{m,k}<x_{n,k}<x_{m,k+1},&\lfloor m/j\rfloor=\lfloor n/j\rfloor+1.
\end{array}
\]
For each fixed $r,k$ there is a constant $C_{r,k}$ such that
\[
 x_{jd+r,k}=-\frac{G}{\lambda_{r,k}}d^j+C_{r,k}d^{j-1}+o(d^{j-1}).
\]
\end{theorem}
The onset may depend on both $j$ and $K$. At $j=2$, the two limiting
functions at $w=-y^2$ are $\cos y$ and $\sin(y)/y$. Along even rows
the limit of $x_{2d,k}/d^2$ is $-32/(9\pi^2(k-1/2)^2)$, and along
odd rows the limit of $x_{2d+1,k}/d^2$ is $-32/(9\pi^2k^2)$.

\section{Reversing the large-block count}
Write $A=\sum_{i=1}^{j-1}z^i/i!$, $B=\sum_{i\ge j}z^i/i!$,
$V=(1-A)^{-1}$, $U=j!B/z^j$ and $H=UV$. The $k$ large blocks
separate $k+1$ possibly empty ordered lists of small blocks. The
labelled product rule gives $[t^k]O_{n,j}/n!=[z^n]B^kV^{k+1}$.
For $n=jd+r$, $d\ge1$, take the polynomial extension at zero of
\[
 R_d(w)=\frac{(j!)^d}{n!(bd)^r}
     \left(\frac{w}{Gd^j}\right)^d O_{n,j}(Gd^j/w).
\]
For $0\le h\le d$, set $m=jh+r$. Then
\[
  a_{d,h}=[w^h]R_d=\frac{[z^m](VH^{d-h})}{(bd)^m}.
\]
Set $a_{d,h}=0$ for $h>d$. The constant terms of $V,H$ are one.
Their linear coefficients
are $1,b$, respectively. Direct multiplication gives
\[
c=[z^2]H=1+\tfrac12{\bf1}_{j\ge3}+\frac1{j+1}+\frac1{(j+1)(j+2)}.
\]

\begin{lemma}\label{corr}
Uniformly on every closed disk, $d(R_d-F_r)\to E_r$, where
$E_r(w)=\sum_{h\ge0}e_hw^h$ and
\[
 e_h=-\frac{mh+m(m-1)/2}{m!}
      +\frac{1/b}{(m-1)!}+\frac{c/b^2}{(m-2)!},\qquad m=jh+r .
\]
A reciprocal factorial with a negative argument denotes zero.
\end{lemma}
\begin{proof}
For fixed $m$ put $D=d-h$. Selecting $m$ linear terms of $H^D$
contributes $b^m\binom Dm$ to $[z^m]VH^D$.
Two other choices have degree $m-1$ in $D$. Taking the linear term of $V$
gives $b^{m-1}\binom D{m-1}$. Taking a quadratic term of $H$ gives
$(m-1)cb^{m-2}\binom D{m-1}$ when $m\ge2$. All further
terms use at most $m-2$ nonconstant factors of $H$. Expansion of
the falling factorial $(D)_m$ proves the coefficient limit, including
$m=1$. For $m=0$ the coefficient is exactly one.

Use $\preceq$ for coefficientwise comparison. Since $A\preceq z/(1-z)$,
the geometric series yields $V\preceq(1-z)/(1-2z)$. Moreover,
$U\preceq1/(1-z)$, so $H\preceq1/(1-2z)$ and both $[z^i]V$ and
$[z^i]H$ are at most $2^i$. For $h\le d$ compare the normalized all-linear
term with $1/m!$ using the factors $\max(0,1-(h+i)/d)$, $0\le i<m$.
Their product vanishes when $D<m$ and agrees with $(D)_m/d^m$ in all cases.
One minus their product is bounded by the sum of their deficits. After
multiplication by $d/m!$ the bound is $(mh+m(m-1)/2)/m!$.
For the other terms the degree drawn from $V$ is positive or some
factor of $H$ contributes at least two. Removing one unit from $V$
or two from a chosen factor of $H$ leaves at most
$\binom{D+m-1}{m-1}$ or $D\binom{D+m-2}{m-2}$ possibilities.
Each term weighs at most $2^m$. Since $m\le jd+r$, their normalized
contribution after multiplication by $d$ is bounded by
\[
 (2/b)^m\left(\frac{(1+j+r)^{m-1}}{(m-1)!}
             +\frac{(1+j+r)^{m-2}}{(m-2)!}\right).
\]
Terms for absent selections are omitted. For $h>d$ the bound
$d/m!\le h/m!$ suffices. These majorants remain summable after
multiplication by $R^h$ for every fixed $R>0$.
\end{proof}

\section{The limiting zeros}
The functions $F_r$ are the Mittag--Leffler functions $E_{j,r+1}$
and the $j$-sections of $e^z$. Popov and Sedletskii~\cite{ps} give
a general treatment of Mittag--Leffler zeros. Athanasiadis and
Wagner~\cite[Theorem 5.3 and Remark 5.4]{aw} establish interlacing
in the weak sense for sections of P\'olya frequency series. We need
simple zeros and a strict, oriented comparison across and within the
sections. Grouping the $j$ consecutive
factorial factors gives $r!F_r={}_0F_{j-1}(\,;{\bf a};w/j^j)$,
where ${\bf a}$ consists of $(r+\ell)/j$, $1\le\ell\le j$, omitting
the entry $1$. Such hypergeometric functions with positive parameters
have only real negative zeros by the Hurwitz--Hille
theorem~\cite[Theorem 1]{sokal}; see also~\cite[Theorem 2]{baricz}.
The same series
for any positive denominator parameters has order $1/j$ by Stirling's
formula. It is nonpolynomial, so Hadamard factorization gives infinitely
many negative zeros and a genus-zero product.

Let $f$ be one such normalized series and $g$ the series obtained by
increasing one parameter $a$ by one. Then $af=ag+wg'$. Away from the
zeros $-u$ of $g$, the derivative on the real axis of
$af/g=a+\sum w/(w+u)$ is $\sum u/(w+u)^2>0$, counting multiplicities.
The logarithmic derivatives converge locally off the poles. Hence a zero
of $f$ away from the zeros of $g$ is simple. At a zero of $g$ of
multiplicity $s$, the multiplicity in $f$ is exactly $s-1$ by
$af=ag+wg'$ and the nonvanishing of $w$. Consequently a multiple zero
of $f$ would persist at the same point through every successive shift of
$a$. The shifted series tend locally uniformly to one since each
nonconstant coefficient decreases to zero and is bounded by its
original value. This proves simplicity.

Coefficient comparison gives $jwF_r'+rF_r=F_{r-1}$ for $1\le r<j$,
and $jF_0'=F_{j-1}$. At $w=-x$ the first ratio is
\[
 \frac{F_{r-1}(-x)}{F_r(-x)}
 =r+j\sum_{\ell\ge1}\frac{x}{x-\lambda_{r,\ell}} .
\]
It decreases strictly from $r$ to $-\infty$ before the first pole,
and from $+\infty$ to $-\infty$ in every following gap.
Simplicity excludes a common zero at a pole. The logarithmic derivative
of $F_0$ has one zero in each interior gap and none before the first pole.
Thus $\lambda_{0,k}<\lambda_{1,k}<\cdots<
\lambda_{j-1,k}<\lambda_{0,k+1}$ for each $k\ge1$.

\section{Transfer to the finite ordered rows}
For a fixed residue take a circle centered at zero with radius between
$\lambda_{r,K+1}$ and $\lambda_{r,K+2}$, and small disjoint circles
about the first $K+1$ limiting zeros inside it. Lemma~\ref{corr}
and Rouch\'e's theorem give one zero counted with multiplicity in each
small disk and exactly $K+1$ inside the large one for all large $d$,
with no zeros on the chosen circles. Conjugation fixes the unique zero
in each small disk. Shrinking the disks shows that the selected roots
$w_{d,k}$ converge to $w_0=-\lambda_{r,k}$. The map $w\mapsto Gd^j/w$
reverses the modulus inequalities. Its derivative is nonzero at these
roots, and the prefactor in $R_d$ is nonzero there, so it preserves
simplicity. The positive constant term of $O_{n,j}$ rules out a root at
zero, and the top coefficient gives $R_d(0)>0$, so all remaining zeros
of the original row have smaller modulus.

The difference quotient of $F_r$ at $w_0$ tends to $F_r'(w_0)\ne0$,
using its continuous extension when $w_{d,k}=w_0$.
Evaluating Lemma~\ref{corr} at $w_{d,k}$ gives
$d(w_{d,k}-w_0)\to-E_r(w_0)/F_r'(w_0)$.
Expansion of $Gd^j/w_{d,k}$ yields Theorem~\ref{thm}'s formula
with $C_{r,k}=GE_r(w_0)/(w_0^2 F_r'(w_0))$.

Write $n=jd+r$ and $m=je+s$. If $e=d$, then $s>r$. Dividing the roots
by $d^j$, the ladder gives $x_{n,k}<x_{m,k}<x_{n,k+1}$ eventually.
Otherwise $e=d+1$ and $s\le r$, since $m-n\le j$. If $s<r$ the ladder
gives $x_{m,k}<x_{n,k}<x_{m,k+1}$ because $(d+1)^j/d^j\to1$.
If $s=r$, then $m=n+j$ and the leading constants at index $k$ coincide.
The common correction coefficient multiplies a difference
$(d+1)^{j-1}-d^{j-1}=o(d^{j-1})$. Thus the two-term expansion gives
\[
 \frac{x_{j(d+1)+r,k}-x_{jd+r,k}}{d^{j-1}}
       \longrightarrow -jG/\lambda_{r,k}<0 .
\]
Different indices are separated at order $d^j$. Taking the maximum
of the onsets for finitely many residues, gaps and indices concludes
the proof.

\paragraph{Code and data availability.}
The accompanying source archive contains two exact arithmetic scripts
using the last-block recurrence and rational coefficient
extraction. The Lean development proves the reciprocal coefficient and
limit formulas, the strict ladder, and the extreme-root window for the
original ordered rows. Its statement map is supplied with the source.

\paragraph{Disclosure.}
ChatGPT assisted with the derivation, mathematical review, formalization
and writing, and the author takes responsibility for the manuscript.

\begin{thebibliography}{9}
\bibitem{catalog} P.~Alexandersson, \emph{Real-rooted polynomials:
open problems}, The Symmetric Functions Catalog, section on set partitions.
\url{https://www.symmetricfunctions.com/realRootedOpenProblems.htm}.
\bibitem{aw} C.~A.~Athanasiadis and D.~G.~Wagner, \emph{Veronese sections
and interlacing matrices of polynomials and formal power series},
\href{https://arxiv.org/abs/2404.12989v2}{arXiv:2404.12989v2} (2026).
\bibitem{baricz} \'A.~Baricz and S.~Singh, \emph{Zeros of some special entire functions},
\href{https://doi.org/10.1090/proc/13927}{Proc. Amer. Math. Soc. 146 (2018), 2207--2216}.
\bibitem{fang} C.~Fang, \emph{Eventual interlacing of large-block
set-partition polynomials},
\href{https://hexagonmath.org/2611.01129v1}{hexagon:2611.01129v1} (2026).
\bibitem{ps} A.~Yu.~Popov and A.~M.~Sedletskii,
\emph{Distribution of roots of Mittag-Leffler functions},
\href{https://doi.org/10.1007/s10958-013-1255-3}{J. Math. Sci. 190 (2013), 209--409}.
\bibitem{sokal} A.~D.~Sokal, \emph{When does a hypergeometric function
\({}_pF_q\) belong to the Laguerre--P\'olya class $LP^+$?},
\href{https://doi.org/10.1016/j.jmaa.2022.126432}{J. Math. Anal. Appl. 515 (2022), 126432}.
\end{thebibliography}
\end{document}
